\chapter{O Sacramento da Penitência e da Reconciliação}

Pelos sacramentos da iniciação cristã, o homem recebe a vida nova de Cristo; todavia, essa vida é levada em ``vasos de barro'' (2 Cor 4, 7), ainda sujeita ao sofrimento, à doença e, sobretudo, ao pecado que pode enfraquecê-la ou até perdê-la. O Senhor Jesus, médico das almas e dos corpos, que perdoou os pecados ao paralítico e lhe restituiu a saúde do corpo, quis que a sua Igreja continuasse, com a força do Espírito Santo, a sua obra de cura e de salvação. É para isso que existem os dois sacramentos de cura: o sacramento da Penitência e o da Unção dos Enfermos. O Catecismo da Igreja Católica apresenta o primeiro desses dois sacramentos nos parágrafos 1422 a 1498, tratando-o como um encontro com a misericórdia divina pelo qual o pecador descobre que a Páscoa de Cristo não é um evento histórico encerrado, mas um poder vivo que continua a operar a reconciliação no coração humano.

\section{Como se chama este sacramento? (1422--1426)}

O Catecismo começa por assinalar que este sacramento recebeu ao longo da tradição vários nomes, cada um deles revelando uma dimensão da riqueza inexaurível do mistério. É chamado \emph{sacramento da conversão} porque realiza sacramentalmente o apelo de Jesus à conversão e o esforço de retornar à casa do Pai, da qual o pecador se afastou pelo pecado. É chamado \emph{sacramento da Penitência} porque consagra uma caminhada pessoal e eclesial de conversão, arrependimento e satisfação por parte do cristão pecador. É chamado \emph{sacramento da confissão} porque o reconhecimento e a confissão dos pecados diante do sacerdote constituem um elemento essencial deste sacramento; e, num sentido mais profundo, é igualmente uma ``confissão'' --- reconhecimento e louvor --- da santidade de Deus e de sua misericórdia para com o homem pecador.

É chamado \emph{sacramento do perdão} porque, pela absolvição sacramental do sacerdote, Deus concede ao penitente ``o perdão e a paz''. É chamado \emph{sacramento da Reconciliação} porque dá ao pecador o amor de Deus que reconcilia: ``Deixai-vos reconciliar com Deus'' (2 Cor 5, 20). Aquele que vive do amor misericordioso de Deus está pronto para responder ao apelo do Senhor: ``Vai primeiro reconciliar-te com teu irmão'' (Mt 5, 24). Esta pluralidade de nomes ilumina facetas complementares de uma mesma realidade: a misericórdia de Deus, que não abandona o filho pródigo, mas vem ao encontro daquele que retorna com o coração contrito.

Por que um sacramento de reconciliação depois do Batismo? O Batismo concede a vida nova de filhos de Deus e perdoa todos os pecados; mas não elimina a fraqueza da natureza humana nem a inclinação ao pecado --- a que a tradição chama concupiscência ---, a qual persiste nos batizados a fim de que prestem suas provas no combate da vida cristã. Assim, ``se dissermos que não temos pecado, enganamo-nos a nós mesmos'' (1 Jo 1, 8). A Igreja sempre soube que era preciso oferecer uma ``segunda tábua de salvação'' para aqueles que, depois do Batismo, caíram em pecado e perderam a graça batismal. O apóstolo João acrescenta: ``Se alguém pecar, temos junto do Pai um advogado: Jesus Cristo, o justo'' (1 Jo 2, 1).

A vida nova recebida na iniciação cristã não suprimiu a fragilidade e a fraqueza da natureza humana, de modo que o combate espiritual continua sendo uma realidade para todo batizado. É nesse combate que o sacramento da Penitência exerce sua função insubstituível: não como concessão à mediocridade, mas como dom de Deus que restaura, fortalece e renova a vida da graça. Os Padres chamavam este sacramento de ``batismo laborioso'', para indicar que, embora não se repita o Batismo, opera-se nele uma conversão análoga --- o retorno ao estado de justificação, pelo qual o pecador é restituído à amizade com Deus e à comunhão com a Igreja.

\section{A conversão dos batizados (1427--1433)}

Jesus chama à conversão. Tal apelo é parte essencial do anúncio do Reino: ``O tempo chegou ao seu termo, o Reino de Deus está próximo: convertei-vos e acreditai na boa-nova'' (Mc 1, 15). O Batismo é o momento principal da primeira e fundamental conversão; mas o apelo de Cristo à conversão não cessa com ele, e continua a fazer-se ouvir na vida dos cristãos. Esta segunda conversão é uma tarefa ininterrupta para toda a Igreja, que contém pecadores no seu seio e que é, ao mesmo tempo, santa e necessitada de purificação, prosseguindo constantemente no seu esforço de penitência e de renovação. O testemunho mais eloquente desta realidade é a conversão de Pedro após a tripla negação: o olhar de misericórdia do Senhor lhe provoca lágrimas de arrependimento e, depois da ressurreição, a tríplice afirmação do seu amor.

Este esforço de conversão não é somente obra humana. É o movimento do ``coração contrito'' (Sl 51, 19) atraído e movido pela graça para responder ao amor misericordioso de Deus, que nos amou primeiro (cf. 1 Jo 4, 10). A penitência interior é uma reorientação radical de toda a vida: retorno a Deus com todo o coração, ruptura com o pecado, aversão ao mal com repugnância pelas más ações cometidas. Ao mesmo tempo, implica o desejo e o propósito de mudar de vida, com a esperança da misericórdia divina e a confiança na ajuda de sua graça. Esta conversão do coração é acompanhada por uma dor e tristeza salutares, a que os Santos Padres chamaram \emph{animi cruciatus} (aflição do espírito) e \emph{compunctio cordis} (compunção do coração).

O coração do homem é pesado e endurecido. É necessário que Deus dê ao homem um coração novo (cf. Ez 36, 26-27). A conversão é, antes de mais, obra da graça de Deus, que faz com que nossos corações se voltem para Ele: ``Convertei-nos, Senhor, e seremos convertidos'' (Lm 5, 21). É ao descobrir a grandeza do amor de Deus que o coração humano é abalado pelo horror e pelo peso do pecado, e começa a temer ofendê-Lo. O coração humano converte-se ao contemplar Aquele a quem os nossos pecados trespassaram, reconhecendo no Crucificado o amor infinito que se oferece como fonte de perdão e de vida nova.

Depois da Páscoa, é o Espírito Santo que ``confunde o mundo no tocante ao pecado'' (Jo 16, 8-9), fazendo ver ao mundo o pecado de não ter acreditado n'Aquele que o Pai enviou. Mas esse mesmo Espírito, que desmascara o pecado, é o Consolador que dá ao coração do homem a graça do arrependimento e da conversão. A segunda conversão tem também uma dimensão comunitária: o apelo se dirige à Igreja inteira --- ``Arrepende-te!'' (Ap 2, 5.16) --- e a comunidade eclesial, ao acolher o pecador que retorna, torna-se instrumento visível da misericórdia divina que nunca se cansa de oferecer recomeços.

\section{As múltiplas formas de penitência na vida cristã (1434--1439)}

A penitência interior do cristão pode ter expressões muito variadas. A Escritura e os Padres insistem sobretudo em três formas clássicas: o \emph{jejum}, a \emph{oração} e a \emph{esmola}, que exprimem a conversão em relação a si mesmo, a Deus e aos outros. A par da purificação radical operada pelo Batismo ou pelo martírio, citam como meios de obter o perdão dos pecados: os esforços realizados para reconciliar-se com o próximo, as lágrimas de penitência, a preocupação com a salvação do próximo, a intercessão dos santos e a prática da caridade que ``cobre uma multidão de pecados'' (1 Pe 4, 8). Todas estas formas expressam, cada uma a seu modo, a disposição interior de quem reconhece a própria fragilidade e se abre à graça.

A conversão realiza-se na vida cotidiana por gestos de reconciliação, pelo cuidado dos pobres, o exercício e a defesa da justiça e do direito, pela confissão das próprias faltas aos irmãos, pela correção fraterna, a revisão de vida, o exame de consciência, a direção espiritual, a aceitação paciente dos sofrimentos e a coragem de suportar a perseguição por amor da justiça. Tomar a sua cruz todos os dias e seguir Jesus é o caminho mais seguro da penitência (cf. Lc 9, 23). Nenhuma dessas práticas é um fim em si mesma; todas elas são pontes que conduzem o coração de volta a Deus, pelas quais o cristão se purifica e cresce na configuração com Cristo.

A Eucaristia e a Penitência guardam entre si uma relação íntima e fecunda. A conversão e a penitência cotidianas têm sua fonte e alimento na Eucaristia, pois nela se torna presente o sacrifício de Cristo que nos reconciliou com Deus; pela Eucaristia nutrem-se e fortalecem-se os que vivem a vida de Cristo, que emprega nela ``o antídoto que nos livra das faltas cotidianas e nos preserva dos pecados mortais''. Os demais meios espirituais --- a leitura da Sagrada Escritura, a oração da Liturgia das Horas e do Pai Nosso, todo ato sincero de culto ou de piedade --- reavivam em nós o espírito de conversão e de penitência, contribuindo para o perdão dos nossos pecados.

Os tempos e os dias de penitência no decorrer do Ano Litúrgico --- a Quaresma, cada sexta-feira em memória da morte do Senhor --- são momentos fortes da prática penitencial da Igreja, particularmente propícios aos exercícios espirituais, às liturgias penitenciais, às peregrinações, ao jejum, à esmola e à partilha fraterna. O Catecismo recorda que a dinâmica da conversão e da penitência foi descrita de modo admirável por Jesus na parábola do filho pródigo, cujo centro é ``o pai misericordioso'': o deslumbramento de uma liberdade ilusória, o abandono da casa paterna, a miséria extrema, o arrependimento, a decisão de declarar-se culpado diante do pai, o caminho do regresso, o acolhimento generoso do pai --- tudo isso são sinais desta vida nova, pura e digna que é a vida do homem que volta para Deus e para o seio da família que é a Igreja.

\section{O sacramento do perdão e da reconciliação (1440--1449)}

O pecado é, antes de mais, ofensa a Deus, ruptura da comunhão com Ele. Ao mesmo tempo, é um atentado contra a comunhão com a Igreja. É por isso que a conversão traz consigo, simultaneamente, o perdão de Deus e a reconciliação com a Igreja, o que é expresso e realizado liturgicamente pelo sacramento da Penitência e da Reconciliação. Somente Deus perdoa os pecados (cf. Mc 2, 7); e Jesus, porque é o Filho de Deus, diz de si próprio: ``O Filho do Homem tem na terra o poder de perdoar os pecados'' (Mc 2, 10). Em virtude de sua autoridade divina, Ele concede este poder aos homens para que o exerçam em seu nome, confiando o exercício do poder de absolvição ao ministério apostólico.

Cristo quis que sua Igreja fosse, toda ela, na sua oração, na sua vida e na sua atividade, sinal e instrumento do perdão e da reconciliação que Ele nos adquiriu ao preço do seu sangue. O apóstolo é enviado ``em nome de Cristo'', e ``é o próprio Deus'' que, por meio dele, exorta e suplica: ``Deixai-vos reconciliar com Deus'' (2 Cor 5, 20). Durante a sua vida pública, Jesus não somente perdoou os pecados, mas manifestou o efeito desse perdão: reintegrou os pecadores perdoados na comunidade do povo de Deus. Sinal claro disto é o fato de Jesus admitir os pecadores à sua mesa e de sentar-se à mesa deles, gesto que exprime ao mesmo tempo, de modo desconcertante, o perdão de Deus e o regresso ao seio do povo de Deus.

Ao tornar os Apóstolos participantes do seu próprio poder de perdoar os pecados, o Senhor deu-lhes também autoridade para reconciliar os pecadores com a Igreja. As palavras ``ligar e desligar'' significam: aquele que vocês excluírem da vossa comunhão ficará também excluído da comunhão com Deus; aquele que de novo receberdes na vossa comunhão, também Deus o acolherá na sua. A reconciliação com a Igreja é inseparável da reconciliação com Deus. Cristo instituiu o sacramento da Penitência para todos os membros pecadores de sua Igreja, especialmente para aqueles que, depois do Batismo, caíram em pecado grave e assim perderam a graça batismal e feriram a comunhão eclesial.

Através das mudanças que a disciplina e a celebração deste sacramento conheceram ao longo dos séculos, distingue-se a mesma estrutura fundamental: por um lado, os atos do homem que se converte sob a ação do Espírito Santo --- contrição, confissão e satisfação ---; por outro, a ação de Deus pela intervenção da Igreja, que, por meio do bispo e seus presbíteros, concede em nome de Jesus Cristo o perdão dos pecados e fixa o modo da satisfação, rezando também pelo pecador e fazendo penitência com ele. A fórmula de absolvição em uso na Igreja latina exprime os elementos essenciais deste sacramento, incluindo uma dimensão de \emph{epiclese}: ``Deus, Pai de misericórdia, que, pela morte e ressurreição de seu Filho, reconciliou o mundo consigo e enviou o Espírito Santo para a remissão dos pecados, te conceda, pelo ministério da Igreja, o perdão e a paz. E eu te absolvo dos teus pecados em nome do Pai, e do Filho e do Espírito Santo.''

\section{Os atos do penitente: contrição, confissão e satisfação (1450--1460)}

Entre os atos do penitente, a contrição ocupa o primeiro lugar. Ela é ``uma dor da alma e uma detestação do pecado cometido, com o propósito de não mais pecar no futuro''. Quando procedente do amor de Deus amado acima de todas as coisas, a contrição é chamada ``perfeita'' (contrição de caridade): ela perdoa as faltas veniais e obtém igualmente o perdão dos pecados mortais, se incluir o propósito firme de recorrer, o quanto antes, à confissão sacramental. A contrição ``imperfeita'' (ou atrição), nascida da consideração da fealdade do pecado ou do temor da condenação eterna, é também um dom de Deus, um impulso do Espírito Santo; ela não obtém por si mesma o perdão dos pecados graves, mas dispõe para obtê-lo no sacramento da Penitência. É conveniente que a recepção deste sacramento seja preparada com um exame de consciência feito à luz da Palavra de Deus.

A confissão dos pecados, mesmo de um ponto de vista simplesmente humano, liberta-nos e facilita a reconciliação com os outros. Pela confissão, o homem enfrenta seus pecados, assume a sua responsabilidade e abre-se de novo a Deus e à comunhão da Igreja, tornando possível um futuro diferente. A confissão ao sacerdote constitui uma parte essencial do sacramento da Penitência: os penitentes devem, na confissão, enumerar todos os pecados mortais de que têm consciência, mesmo que sejam secretíssimos. Sem ser estritamente necessária, a confissão das faltas cotidianas (pecados veniais) é, contudo, vivamente recomendada pela Igreja, pois ajuda a formar a consciência, a lutar contra as más inclinações, a deixar-se curar por Cristo e a progredir na vida do Espírito.

Muitos pecados prejudicam o próximo. Há que fazer o possível para reparar o dano causado: restituir coisas roubadas, restabelecer a boa reputação de quem foi caluniado, indenizar por ferimentos. A simples justiça o exige. Mas, além disso, o pecado fere e enfraquece o próprio pecador e suas relações com Deus e com o próximo. A absolvição tira o pecado, mas não remedeia todas as desordens por ele causadas. Aliviado do pecado, o penitente deve ainda recuperar a perfeita saúde espiritual, ``satisfazendo'' de modo apropriado ou ``expiando'' seus pecados. A esta satisfação também se dá o nome de ``penitência''.

A penitência que o confessor impõe deve ter em conta a situação pessoal do penitente e procurar seu bem espiritual. Deve corresponder, o quanto possível, à gravidade e à natureza dos pecados cometidos. Pode consistir na oração, em uma doação, nas obras de misericórdia, no serviço ao próximo, em privações voluntárias, sacrifícios e, sobretudo, na aceitação paciente da cruz que temos de carregar. Tais penitências ajudam-nos a configurar-nos com Cristo, que, por si só, expiou os nossos pecados de uma vez por todas (cf. Rm 3, 25; 1 Jo 2, 1-2), e fazem que nos tornemos co-herdeiros de Cristo Ressuscitado, ``uma vez que também sofremos com Ele'' (Rm 8, 17).

\section{O ministro do sacramento (1461--1467)}

Uma vez que Cristo confiou aos Apóstolos o ministério da reconciliação, os bispos, seus sucessores, e os presbíteros, colaboradores dos bispos, continuam a exercer tal ministério. Com efeito, os bispos e os presbíteros têm, em virtude do sacramento da Ordem, o poder de perdoar todos os pecados em nome do Pai e do Filho e do Espírito Santo. O perdão dos pecados reconcilia com Deus, mas também com a Igreja; por isso, o bispo, chefe visível da Igreja particular, é considerado, desde os tempos antigos, o principal detentor do poder e ministério da reconciliação e o moderador da disciplina penitencial, enquanto os presbíteros o exercem na medida em que receberam o encargo do seu bispo ou do Papa.

Certos pecados particularmente graves são punidos com a excomunhão, a pena eclesiástica mais severa, que impede a recepção dos sacramentos e o exercício de certos atos eclesiásticos; a absolvição desses pecados só pode ser dada, segundo o direito da Igreja, pelo Papa, pelo bispo do lugar ou por sacerdotes por eles autorizados. Em caso de perigo de morte, porém, qualquer sacerdote, mesmo que careça da faculdade de ouvir confissões, pode absolver de qualquer pecado e de toda excomunhão. Os sacerdotes devem exortar os fiéis a aproximarem-se do sacramento da Penitência e mostrar-se disponíveis para a sua celebração sempre que os cristãos o peçam de modo razoável.

Ao celebrar o sacramento da Penitência, o sacerdote exerce o ministério do bom Pastor que procura a ovelha perdida, do bom Samaritano que cura as feridas, do pai que espera pelo filho pródigo e o acolhe no seu regresso, do justo juiz cujo julgamento é, ao mesmo tempo, justo e misericordioso. Em resumo, o sacerdote é sinal e instrumento do amor misericordioso de Deus para com o pecador. O confessor não é senhor, mas servidor do perdão de Deus. O ministro deste sacramento deve unir-se à intenção e à caridade de Cristo; deve possuir um conhecimento provado do comportamento cristão, experiência das coisas humanas, respeito e delicadeza para com aquele que caiu, amor à verdade e fidelidade ao Magistério da Igreja.

Dada a delicadeza e a grandeza deste ministério e o respeito devido às pessoas, a Igreja declara que todo sacerdote que ouve confissões está obrigado a guardar segredo absoluto sobre os pecados que seus penitentes lhe confessaram, sob penas gravíssimas. Tampouco pode servir-se dos conhecimentos que a confissão lhe proporciona sobre a vida dos penitentes. Este segredo, que não admite exceções, é chamado ``sigilo sacramental'', porque aquilo que o penitente manifestou ao sacerdote fica ``selado'' pelo sacramento. Não existe nenhuma causa, por mais grave que seja, que possa justificar a violação do sigilo sacramental; ele é, portanto, inteiramente absoluto.

\section{Os efeitos do sacramento (1468--1470)}

``Toda a eficácia da Penitência consiste em nos restituir à graça de Deus e em nos unir a Ele numa amizade perfeita.'' O fim e o efeito deste sacramento são, portanto, a reconciliação com Deus. Naqueles que recebem o sacramento da Penitência com coração contrito e disposição religiosa, seguem-se ``a paz e a tranquilidade da consciência, acompanhadas de uma grande consolação espiritual''. Com efeito, o sacramento da reconciliação com Deus conduz a uma verdadeira ``ressurreição espiritual'': à restituição da dignidade e dos bens próprios da vida dos filhos de Deus, o mais precioso dos quais é a amizade do mesmo Deus. Esta ressurreição espiritual não é meramente metafórica; é uma transformação real pela qual a vida da graça, morta pelo pecado, é genuinamente restaurada.

Este sacramento reconcilia-nos também com a Igreja. O pecado abala ou rompe a comunhão fraterna; o sacramento da Penitência repara ou restaura essa comunhão. Nesse sentido, ele não se limita a curar aquele que é restabelecido na comunhão eclesial, mas também exerce um efeito vivificante sobre a vida da Igreja que sofreu com o pecado de um de seus membros. Restabelecido ou confirmado na comunhão dos santos, o pecador é fortalecido pela troca de bens espirituais entre todos os membros vivos do corpo de Cristo. Como afirmou João Paulo~II: ``A reconciliação com Deus tem como consequência outras reconciliações: o penitente reconcilia-se consigo mesmo no mais profundo do seu ser [...] reconcilia-se com os irmãos; reconcilia-se com a Igreja; reconcilia-se com toda a criação.''

Neste sacramento, o pecador, remetendo-se ao juízo misericordioso de Deus, antecipa de certo modo o julgamento a que será submetido no fim desta vida terrena. É aqui e agora, nesta vida, que nos é oferecida a opção entre a vida e a morte. Só pelo caminho da conversão podemos entrar no Reino de que o pecado grave nos exclui. Convertendo-se a Cristo pela penitência e pela fé, o pecador passa da morte à vida ``e não é submetido a julgamento'' (Jo 5, 24). Este efeito escatológico da Penitência revela que este sacramento não é somente uma correção moral, mas uma verdadeira passagem --- da morte para a vida, do pecado para a graça, do isolamento para a comunhão.

Os efeitos do sacramento da Penitência são, portanto, múltiplos e profundos: reconciliação com Deus pela qual o penitente recupera a graça; reconciliação com a Igreja; remissão da pena eterna em que incorreu pelos pecados mortais; remissão, ao menos em parte, das penas temporais consequentes ao pecado; paz e serenidade da consciência; consolação espiritual; e acréscimo das forças espirituais para o combate cristão. Estes efeitos dependem não de uma afirmação psicológica subjetiva, mas da ação objetiva de Cristo que opera pelo ministério do sacerdote; por isso, são reais e garantidos para quem se aproxima do sacramento com as devidas disposições.

\section{As indulgências (1471--1479)}

A doutrina e a prática das indulgências na Igreja estão estreitamente ligadas aos efeitos do sacramento da Penitência. A definição clássica é a seguinte: ``A indulgência é a remissão, perante Deus, da pena temporal devida aos pecados cuja culpa já foi apagada; remissão que o fiel devidamente disposto obtém em certas e determinadas condições, pela ação da Igreja, a qual, enquanto dispensadora da redenção, distribui e aplica por sua autoridade o tesouro das satisfações de Cristo e dos santos.'' A indulgência é parcial ou plenária, conforme liberta parcialmente ou na totalidade da pena temporal devida ao pecado; e o fiel pode obtê-la para si mesmo ou aplicá-la pelos defuntos.

Para compreender esta doutrina, é preciso ter presente que o pecado tem uma dupla consequência. O pecado grave priva-nos da comunhão com Deus e, portanto, torna-nos incapazes da vida eterna: esta privação chama-se ``pena eterna'' do pecado. Por outro lado, todo pecado, mesmo venial, traz consigo um apego desordenado às criaturas que precisa ser purificado, seja nesta vida, seja depois da morte no estado chamado Purgatório: esta purificação liberta do que se chama ``pena temporal'' do pecado. O perdão do pecado e o restabelecimento da comunhão com Deus trazem consigo a abolição das penas eternas; mas as penas temporais podem subsistir, e o cristão deve esforçar-se por acolhê-las como graça, suportando com paciência os sofrimentos e as provações.

O cristão que busca purificar-se do pecado não se encontra sozinho. A vida de cada filho de Deus está ligada de modo admirável, em Cristo e por Cristo, à vida de todos os outros irmãos cristãos, ``na unidade sobrenatural do Corpo Místico de Cristo, como que numa pessoa mística''. Na comunhão dos santos, existe entre os fiéis --- os que já estão na pátria celeste, os admitidos à expiação do Purgatório e os que ainda vivem peregrinando na terra --- um constante vínculo de amor e uma abundante troca de todos os bens. A santidade de um aproveita aos demais, bem além do dano que o pecado de um possa ter causado aos outros. A esses bens espirituais da comunhão dos santos também se chama ``o tesouro da Igreja''.

Este tesouro não é um somatório de bens materiais acumulados no decurso dos séculos, mas o preço infinito e inesgotável que têm junto de Deus as expiações e méritos de Cristo, nosso Senhor, oferecidos para que a humanidade seja liberta do pecado e chegue à comunhão com o Pai. A este tesouro pertencem também as orações e boas obras da Bem-aventurada Virgem Maria e de todos os santos. A indulgência é obtida mediante a Igreja que, em virtude do poder de ligar e desligar que lhe foi concedido por Jesus Cristo, intervém a favor de um cristão e abre para ele o tesouro dos méritos de Cristo e dos santos. Uma vez que os fiéis defuntos em vias de purificação também são membros da mesma comunhão dos santos, nós podemos ajudá-los, entre outros modos, obtendo para eles indulgências, a fim de que sejam libertados das penas temporais devidas pelos seus pecados.

\section{A celebração do sacramento da Penitência (1480--1484)}

Tal como todos os sacramentos, a Penitência é uma ação litúrgica. Ordinariamente, os elementos de sua celebração são os seguintes: a saudação e a bênção do sacerdote; a leitura da Palavra de Deus para iluminar a consciência e suscitar a contrição; a exortação ao arrependimento; a confissão que reconhece os pecados e os manifesta ao sacerdote; a imposição e a aceitação da penitência; a absolvição do sacerdote; e o louvor de ação de graças com a despedida e a bênção. Cada um desses elementos tem sua significação própria e contribui para a integridade do ato sacramental que opera a reconciliação.

A liturgia bizantina tem várias fórmulas de absolvição em forma deprecativa, que exprimem admiravelmente o mistério do perdão: ``Deus, que pelo profeta Natã perdoou a Davi, quando ele confessou os seus próprios pecados, a Pedro depois de ele ter chorado amargamente [...] este mesmo Deus vos perdoe, por intermédio de mim pecador, nesta vida e na outra, e vos faça comparecer, sem vos condenar, no seu temível tribunal.'' O sacramento da Penitência pode também ter lugar no âmbito de uma celebração comunitária, na qual se faz uma preparação conjunta para a confissão e conjuntamente se dão graças pelo perdão recebido; neste caso, a confissão pessoal dos pecados e a absolvição individual são inseridas em uma liturgia da Palavra com leituras, homilia e exame de consciência feito em comum, pedido comunitário de perdão e Oração do Senhor.

Em casos de grave necessidade, pode-se recorrer à celebração comunitária com confissão geral e absolvição geral. Tal necessidade pode ocorrer quando há perigo iminente de morte, sem que o sacerdote ou os sacerdotes tenham tempo suficiente para ouvir a confissão de cada penitente; ou quando, tendo em conta o número dos penitentes, não há confessores bastantes para ouvir devidamente as confissões individuais em tempo razoável. Para a validade da absolvição nesses casos, os fiéis devem ter o propósito de confessar individualmente seus pecados graves em tempo oportuno. Pertence ao bispo diocesano julgar se as condições requeridas para a absolvição geral estão presentes.

``A confissão individual e íntegra e a absolvição constituem o único modo ordinário pelo qual o fiel, consciente de pecado grave, se reconcilia com Deus e com a Igreja: somente a impossibilidade física ou moral o escusa desta forma de confissão.'' Há razões profundas para que assim seja: Cristo age em cada um dos sacramentos de modo pessoal, dirigindo-se a cada pecador --- ``Meu filho, os teus pecados são-te perdoados'' (Mc 2, 5) --- como o médico que se inclina sobre cada doente. A confissão pessoal é, portanto, a forma mais significativa da reconciliação com Deus e com a Igreja, aquela em que melhor se manifesta a ternura pessoal de Cristo para com cada alma.

\section{Em síntese (1485--1498)}

``Na tarde da Páscoa, o Senhor Jesus apareceu aos seus Apóstolos e disse-lhes: "Recebei o Espírito Santo: àqueles a quem perdoardes os pecados, ser-lhes-ão perdoados; e àqueles a quem os retiverdes, ser-lhes-ão retidos"'' (Jo 20, 22-23). Por esta palavra solene, Cristo institui o sacramento que a Igreja chama sacramento da conversão, da confissão, da Penitência ou da Reconciliação. O perdão dos pecados cometidos depois do Batismo é concedido por meio desse sacramento próprio: quem peca ofende a honra de Deus e seu amor, a própria dignidade de homem chamado a ser filho de Deus, e o bem-estar espiritual da Igreja, da qual cada fiel deve ser pedra viva. Aos olhos da fé, não existe mal mais grave do que o pecado; nada tem piores consequências para os próprios pecadores, para a Igreja e para todo o mundo.

Voltar à comunhão com Deus, depois de a ter perdido pelo pecado, é um movimento nascido da graça do mesmo Deus misericordioso e cheio de interesse pela salvação dos homens. O movimento de regresso a Deus, pela conversão e arrependimento, implica dor e aversão em relação aos pecados cometidos e o propósito firme de não tornar a pecar no futuro. A conversão refere-se ao passado e ao futuro; alimenta-se da esperança na misericórdia divina. O sacramento da Penitência é constituído pelo conjunto de três atos realizados pelo penitente --- arrependimento, confissão ou manifestação dos pecados ao sacerdote, e propósito de cumprir a reparação --- e pela absolvição do sacerdote. O arrependimento, se motivado pelo amor de caridade para com Deus, é ``perfeito''; se fundado em outros motivos, é ``imperfeito''.

Aquele que quer obter a reconciliação com Deus e com a Igreja deve confessar ao sacerdote todos os pecados graves que ainda não tiver confessado e de que se lembre depois de ter examinado cuidadosamente a consciência. A confissão das faltas veniais, sem ser em si necessária, é todavia vivamente recomendada pela Igreja. O confessor propõe ao penitente o cumprimento de certos atos de ``satisfação'' ou ``penitência'', com o fim de reparar o mal causado pelo pecado e restabelecer os hábitos próprios de um discípulo de Cristo. Somente os sacerdotes que receberam da autoridade da Igreja a faculdade de absolver podem perdoar os pecados em nome de Cristo.

Os efeitos espirituais do sacramento da Penitência são: a reconciliação com Deus, pela qual o penitente recupera a graça; a reconciliação com a Igreja; a remissão da pena eterna em que incorreu pelos pecados mortais; a remissão, ao menos em parte, das penas temporais consequentes ao pecado; a paz e a serenidade da consciência e a consolação espiritual; e o acréscimo das forças espirituais para o combate cristão. A confissão individual e integral dos pecados graves, seguida da absolvição, continua a ser o único meio ordinário para a reconciliação com Deus e com a Igreja. Por meio das indulgências, os fiéis podem obter para si mesmos, e também para as almas do Purgatório, a remissão das penas temporais consequentes ao pecado.
